\begin{lemma}[Кронекера (б/д)]
Пусть \( \sum \limits_{n = 1}^{\infty} x_n < \infty , 0 < b_n \uparrow \infty \).
Тогда \( \frac{\sum \limits_{k = 1}^n x_k b_k}{b_n} \to 0, n \to \infty \).
\end{lemma}

\begin{theorem}[УЗБЧ 2]
Пусть \( \xi_1, \xi_2, \dots \) -- н. о. р. с. в. (независимые одинаково распределённые
случаные величины), \( \E |\xi_1| < \infty \). Тогда
\( \frac{S_n - \E S_n}{n} \xrightarrow{\text{ п. н. }} 0 \).
\end{theorem}

\begin{proof}
Без ограничения общности будем считать, что \( \E \xi_1  = 0\). 

Так как \( \E |\xi_1| < \infty \), то можно написать такую оценку:

\[
\sum_{k = 1}^{\infty} (k - 1) \cdot I(k - 1 \leqslant |\xi_1| < k) 
\leqslant |\xi_1| \leqslant
\sum_{k = 1}^{\infty} k \cdot I(k - 1 \leqslant |\xi_1| < k)
\]

Здесь просто на каждом промежутке \( [k - 1, k) \) ограничили
\( \xi_1 \) снизу как \( k - 1 \), сверху -- как \( k \).
Теперь заметим, что этот ряд можно переписать следующим образом:

\begin{multline*}
\sum_{k = 1}^{\infty} k \cdot I(k - 1 \leqslant |\xi_1| < k) = 
1 \cdot I(0 \leqslant |\xi_1| < 1) + 
2 \cdot I(1 \leqslant |\xi_1| < 2) +
3 \cdot I(2 \leqslant |\xi_1| < 3) +
\dots = \\
= I(|\xi_1| \geqslant 0) + 1 \cdot I(1 \leqslant |\xi_1| < 2) +
2 \cdot I(2 \leqslant |\xi_1| < 3) + \dots = 
\sum_{k = 1}^{\infty} I(|\xi_1| \geqslant k - 1)
\end{multline*}

Аналогично переписав другой ряд, получаем оценку:

\[
\sum_{k = 1}^{\infty} I(|\xi_1| \geqslant k) \leqslant |\xi_1| 
\leqslant \sum_{k = 1}^{\infty} I(|\xi_i| \geqslant k - 1)
\]

Тогда матожидание оценивается следующим образом:

\[
\sum_{k = 1}^{\infty} P(|\xi_1| \geqslant k) \leqslant \E |\xi_1|
\leqslant \sum_{k = 1}^{\infty} P(|\xi_1| \geqslant k - 1) = 
1 + \sum_{k = 1}^{\infty} P(|\xi_1| \geqslant k)
\]

Верхняя и нижняя оценки отличаются ровно на единицу, поэтому 
\( \E |\xi_1| < \infty \) тогда и только тогда, когда
\( \sum \limits_{k = 1}^{\infty} P(|\xi_1| \geqslant k) < \infty \). 
В силу одинаковой распределённости случайных величин последнее условие
можно заменить на 
\( \sum \limits_{k = 1}^{\infty} P(|\xi_k| \geqslant k) < \infty \).
По лемме Борелля-Кантелли сходимость этого ряда означает, что
\( P( |\xi_n| \geqslant n \text{ б. ч. }) = 0 \). То есть начиная с
некоторого момента все $\xi_n$ будут меньше чем $n$. Тогда будет логично
рассматривать случайные величины 
\( \widetilde{\xi_n} = \xi_n I(|\xi_n| \leqslant n) \).
Пусть \( A = \overline{\{ |\xi_n| \geqslant n \text{ б. ч. } \}}\).

Тогда \( \forall \omega \in A \;\; \exists n_0 : \forall n \geqslant n_0 \;\;
\widetilde{\xi_n} = \xi_n \). Из этого условия и того, что \( P(A) = 1 \),
можно сделать вывод, что
на пределы 

\[
\lim_{n \to \infty} \frac{\xi_1 + \dots + \xi_n}{n} \text{ и }
\lim_{n \to \infty} \frac{\widetilde{\xi_1} + \dots + \widetilde{\xi_n}}{n}
\]

одновременно существуют и равны на \( A \) (*).

В силу одинаковой распределённости случайных величин получаем, что
\( \E \widetilde{\xi_n} = \E \xi_n I(|\xi_n| \leqslant n) = 
\E \xi_1 I(|\xi_1| \leqslant n) \). Заметим, что
\( \xi_1 I(|\xi_1| \leqslant n) \xrightarrow{\text{ п. н. }} \xi_1 \).
Тогда можем применить теорему Лебега о мажорируемой сходимости 
(выражение \( \xi_1 I(|\xi_1| \leqslant n) \) ограничивается \( \xi_1 \),
у которой конечное матожидание): 

\[
\E \widetilde{\xi_n} = \E \xi_1 I(|\xi_1| \leqslant n) \to 
\E \xi_1 = 0
\]

Тогда последовательность сходится и в смысле средних арифметических:

\[
\frac{\E \widetilde{\xi_ 1} + \dots + \E \widetilde{\xi_n}}{n} \to 0 \; (**)
\]

Теперь из (*) и (**) можно сделать вывод, что 

\[
P \left[ \lim_{n \to \infty} \frac{\xi_1 + \dots + \xi_n}{n} =
\lim_{n \to \infty} \frac{\widetilde{\xi_1} + \dots + \widetilde{\xi_n}
- \E \widetilde{\xi_1} - \dots - \E \widetilde{\xi_n}
}{n}
\right] = 1
\]

Если мы докажем, что 
\(\sum \limits_{k = 1}^{\infty} \frac{\D \widetilde{\xi_k}}{k^2} < \infty\),
то из этого будет следовать (по теореме Колмогорова-Хинчина), что
\( \sum \limits_{k = 1}^{\infty} 
\frac{\widetilde{\xi_k} - \E \widetilde{\xi_k}}{k} < \infty \).
Отсюда, в свою очередь, по лемме Кронекера будет следовать, что
\(
\frac{1}{n} \sum \limits_{k = 1}^n 
\frac{\widetilde{\xi_k} - \E \widetilde{\xi_k}}{k} \cdot k = 
\frac{\widetilde{\xi_1} + \dots + \widetilde{\xi_n} - 
\E \widetilde{\xi_1} - \dots - \E \widetilde{\xi_n}}{n}
\xrightarrow{\text{ п. н. }} 0
\). А значит, по доказанному, и 
\(
\frac{\xi_1 + \dots + \xi_n}{n} \xrightarrow{\text{ п. н. }} 0
\)

Осталось доказать, что 
\( \sum \limits_{k = 1}^{\infty} 
\frac{\D \widetilde{\xi_k}}{k^2} < \infty \).

\begin{multline*}
\sum_{k = 1}^{\infty} \frac{\D \widetilde{\xi_k}}{k^2} \leqslant
\sum_{k = 1}^{\infty} \frac{\E \widetilde{\xi_k}^2}{k^2} = 
\sum_{k = 1}^{\infty} \frac{\E \xi_k^2 I(|\xi_k| \leqslant k)}{k^2} = 
\sum_{k = 1}^{\infty} \frac{\E \xi_1^2 I(|\xi_1| \leqslant k)}{k^2} = 
\sum_{k = 1}^{\infty} \frac{1}{k^2} \sum_{j = 1}^k 
\E \xi_1^2 I(j - 1 < |\xi_1| \leqslant j) = \\
= \sum_{j = 1}^{\infty} \sum_{k = j}^{\infty} \frac{1}{k^2}
\E \xi_1^2 I(j - 1 < |\xi_1| \leqslant j) = 
\sum_{j = 1}^{\infty} \E \xi_1^2 I(j - 1 < |\xi_1| \leqslant j)
\sum_{k = j}^{\infty} \frac{1}{k^2} \leqslant
\sum_{j = 1}^{\infty} \E \xi_1^2 I(j - 1 < |\xi_1| \leqslant j) \cdot
\frac{2}{j} \leqslant \\
\leqslant \sum_{j = 1}^{\infty} j^2 P(j - 1 < |\xi_1| \leqslant j) \cdot
\frac{2}{j} = 
2 \sum_{j = 1}^{\infty} j P(j - 1 < |\xi_1| \leqslant j) = 
2 \sum_{j = 1}^{\infty} (j - 1) P(j - 1 < |\xi_1| \leqslant j) + \\
+ 2 \sum_{j = 1}^{\infty} P(j - 1 < |\xi_1| \leqslant j) \leqslant
2 \E |\xi_1| + 2 < \infty
\end{multline*}

\end{proof}